\subsection{Statistical Significance Analysis}

Reporting mean performance alone does not establish whether observed differences between models are attributable to chance, particularly when evaluation is based on a limited number of cross-validation folds. We therefore conduct formal statistical validation to support claims of superiority over competing approaches.

We employ the Wilcoxon signed-rank test, a non-parametric paired test that does not assume normally distributed fold-wise scores and is well suited to small sample sizes typical of $k$-fold cross-validation. Comparisons are conducted between the proposed contrastive self-supervised model and each major baseline (ResNet-50, ViT, EfficientNet, ConvNeXt, and AST) using fold-wise F1-scores collected from repeated 5-fold cross-validation ($n = 5 \times 2$ paired measurements per comparison). Statistical significance is assessed at $\alpha = 0.05$ (two-sided).

Table~\ref{tab:wilcoxon} summarizes the results. The proposed model attains a mean F1 of 92.37\% $\pm$ 0.29\% across folds. All pairwise comparisons yielded $p < 0.05$, indicating that the proposed model achieves statistically significant improvements in F1-score over every competing baseline under repeated 5-fold cross-validation.

\begin{table}[htbp]
\centering
\caption{Statistical significance of the proposed model versus competing baselines (Wilcoxon signed-rank test on fold-wise F1; proposed model: 92.37\% $\pm$ 0.29\%).}
\label{tab:wilcoxon}
\begin{tabular}{lccc}
\hline
Baseline Model & Mean $\pm$ Std (\%) & $p$-value & Significant \\
\hline
AST & 89.92 $\pm$ 0.29 & 0.0020 & Yes \\
ConvNeXt & 88.05 $\pm$ 0.23 & 0.0020 & Yes \\
EfficientNet & 87.24 $\pm$ 0.40 & 0.0020 & Yes \\
ViT & 86.12 $\pm$ 0.18 & 0.0020 & Yes \\
ResNet-50 & 84.47 $\pm$ 0.36 & 0.0020 & Yes \\
\hline
\end{tabular}
\end{table}


